% TOP_PROBLEM: {"id":30007560,"problem_number":"LOCAL-30007560","title":"Fair allocation across changing profiles","category_id":21,"category":{"id":21,"name":"economics","display_name":"Economics","description":"Open economic theory statements, published research agendas, and proposed scoped specifications; question provenance and historical priority are qualified per record.","slug":"economics","order_index":21,"created_at":"2026-10-08T00:00:00Z"},"status":"research_question","external_url":null,"legacy_ids":[],"published":false,"statement":"Characterize which combinations of EF1, Pareto efficiency, resource monotonicity, and population monotonicity can coexist on the specified finite additive-valuation domains.","economics":{"original_id":49,"field":"Economic theory, equilibrium and social choice","kind":"theoretical","formalization_basis":"adapted_research_question","source_status":"research_agenda","priority_status":"earliest_found","priority_note":"This is the earliest source in which the relevant agenda was verified during this audit; absolute priority has not been established.","earliest_verified_reference":{"authors":["Dominik Peters"],"title":"Inter-profile Axioms in the Allocation of Indivisible Goods, in Fair Division: Algorithms, Solution Concepts, and Applications","year":2024,"url":"https://drops.dagstuhl.de/storage/04dagstuhl-reports/volume14/issue09/24401/html/DagRep.14.9.145/DagRep.14.9.145.html","doi":"10.4230/DagRep.14.9.145","locator":"Seminar report §4.5","evidence_summary":"The contribution asks for allocation rules or impossibility results combining fairness with comparisons across resource, population, and valuation profiles."},"current_reference":{"authors":["Dominik Peters"],"title":"Inter-profile Axioms in the Allocation of Indivisible Goods, in Fair Division: Algorithms, Solution Concepts, and Applications","year":2024,"url":"https://drops.dagstuhl.de/storage/04dagstuhl-reports/volume14/issue09/24401/html/DagRep.14.9.145/DagRep.14.9.145.html","doi":"10.4230/DagRep.14.9.145","locator":"Seminar report §4.5","evidence_summary":"The contribution asks for allocation rules or impossibility results combining fairness with comparisons across resource, population, and valuation profiles."},"formalization_note":"The finite-domain compatibility problem is an adaptation of the stated inter-profile agenda. Known impossibilities must be excluded or recorded during specialist review."},"statusQualification":"Published question or research agenda verified at the cited locator. The mathematical specification is adapted for this catalogue and includes additional assumptions; source publication does not establish first-ever priority or certify this exact model remains unsolved. The finite-domain compatibility problem is an adaptation of the stated inter-profile agenda. Known impossibilities must be excluded or recorded during specialist review.","statusReviewedAt":"2026-10-08"}
% FIRST_POSTED: 2026-10-08
% ENTRY_KIND: statement_only
% !TeX program = lualatex
\documentclass[11pt]{article}
\usepackage{fontspec}
\usepackage[a4paper,margin=25mm]{geometry}
\usepackage{amsmath,amssymb,mathtools}
\usepackage{xurl}
\usepackage[hidelinks]{hyperref}
\newcommand{\sref}[2]{\hyperlink{source:#1:#2}{[S#2]}}
\setlength{\emergencystretch}{3em}
\begin{document}
% problemId:30007560
\section*{Fair allocation across changing profiles}\label{30007560}
\subsection{Definitions and mathematical statement}
Fix a finite universe of agents and goods and integer \(V\geq1\). Each profile consists of agent set \(N\), good set \(G\), and additive valuations \(v_i(g)\in\{0,\ldots,V\}\). A deterministic allocation rule \(F\) selects a complete allocation for every profile in a specified domain \(\mathcal D\). It is envy-free up to one good if, for each ordered pair \(i,j\), either \(F_j=\varnothing\) or some \(g\in F_j\) satisfies \(v_i(F_i)\geq v_i(F_j\setminus\{g\})\). It is Pareto efficient if no other complete allocation weakly raises every agent's value and strictly raises at least one. Resource monotonicity requires that adding a good without changing old valuations weakly increases every incumbent's value. Population monotonicity requires that adding an agent without changing incumbent valuations weakly decreases every incumbent's value. Determine which subsets of these four requirements are jointly attainable on \(\mathcal D\), and characterize the maximal domains or explicit weakenings that permit compatibility. Quantifiers must range over every admissible profile and each permitted profile transition, rather than only selected examples. Supply constructive rules or impossibility proofs, with the relevant valuation and population restrictions made explicit. This asks for a compatibility frontier; it does not conjecture that the strongest bundle of axioms is feasible on unrestricted profiles.

\par\textbf{Formulation and scope.} The finite-domain compatibility problem is an adaptation of the stated inter-profile agenda. Known impossibilities must be excluded or recorded during specialist review.

\par\textbf{Open-question provenance.} An explicit question or research agenda was verified at the source locator. This does not by itself certify that every later version remains unresolved \sref{30007560}{1}.

\par\textbf{Historical priority.} This is the earliest source in which the relevant agenda was verified during this audit; absolute priority has not been established.

\subsection{Short English statement}
Characterize which combinations of EF1, Pareto efficiency, resource monotonicity, and population monotonicity can coexist on the specified finite additive-valuation domains.

\subsection{Sources}
\begin{itemize}\raggedright
\item[\textnormal{[S1]}]\hypertarget{source:30007560:1}{} Dominik Peters. 2024. Inter-profile Axioms in the Allocation of Indivisible Goods, in Fair Division: Algorithms, Solution Concepts, and Applications. Seminar report §4.5.
\url{https://drops.dagstuhl.de/storage/04dagstuhl-reports/volume14/issue09/24401/html/DagRep.14.9.145/DagRep.14.9.145.html}.
DOI: 10.4230/DagRep.14.9.145.
{\small\textit{Earliest verified question/agenda reference; historical priority qualified below. The contribution asks for allocation rules or impossibility results combining fairness with comparisons across resource, population, and valuation profiles.}}
\end{itemize}
\end{document}
